\section{Conclusion}

Forecasts made during accelerating technological change fail when they preserve stale execution assumptions or translate capability directly into realized value. The proposed framework permits acceleration and discontinuity, but forces every projection through evidence comparability, verification, authority, adoption, demand, and residual constraints. It recomputes work from task graphs, compares competing regimes under a common likelihood, propagates uncertainty through full scenario paths, exposes measurement breaks, and converts monitoring into dated action rules.

The compact canonical model is intentionally simple enough to remember:
\[
Y_s(t)=B_s(t)\min\{\text{technical capacity},\text{absorptive demand}\}.
\]
Its disciplined use is not simple. It requires fresh evidence, compatible units, explicit elasticities, a double-counting audit, first-principles execution analysis, calibrated humility about regime transitions, and repeated recalculation. The generalized realization family extends the model only when branch-specific constraints or partial substitution matter enough to justify the added identification burden.

Version 2.0.0 turns the method into a coordinated research system: a complete executable prompt, strict schemas, an end-to-end engine, probabilistic hindcasts, proper scoring, a prospective hash chain, an SBOM, reproducible release machinery, and explicit external validation gates. The package is designed to be falsified and improved, not insulated from evidence.

The central practical instruction remains:

\begin{keybox}
\centering\large\bfseries
Never extrapolate stale timelines. Forecast a world that may accelerate during the forecast itself.
\end{keybox}
